\worldchapter{\trans{foes_title}}{appendix-foes}
\label{ch:appendix-foes}

\begin{multicols}{2}


    \section{Untoter}



    \subsection*{Zombie}

    Der Zombie ist eine tragische Gestalt, welche in vielen Geschichten vorkommt. Ein Mensch, welcher auf besondere und unnatürliche Weise nach seinem Tod am Leben gehalten wird. Zombies haben ein Gehirn von der Größe einer Erbse, sie kennen nicht viel mehr als das Verlangen nach Blut. Und Gehirnen. Also, wenn man der Popkultur glauben darf.

    \begin{tabular}{@{}ll@{}}
        \textbf{\trans{wounds_label}:} & 6 \\
        \textbf{\trans{movement_label}:} & 3 \\
        \textbf{\trans{strength_label}:} & 3 \\
        \textbf{\trans{dexterity_label}:} & 1 \\
        \textbf{\trans{mind_label}:} & 1 \\
        \textbf{\trans{resistances_label}:} & ['Feuer', 'Gift'] \\
    \end{tabular}

    \textbf{Untoter Griff} (4)\\
    Vergiftet 1\\
    \textbf{Biss} (5)\\
    Vergiftet 1\\


    \subsection*{Mumie}

    Die lebendige Mumie ist ein Untoter. Sie ist auf magische Weise zum Leben erweckt worden. Nichts ist von ihrem Geist erhalten geblieben, alles, wonach sie trachtet, ist das Leben ihrer Opfer zu nehmen. Sie ist in der Regel unbewaffnet, ihr Fluch vergiftet jedoch ihre Opfer.

    \begin{tabular}{@{}ll@{}}
        \textbf{\trans{wounds_label}:} & 8 \\
        \textbf{\trans{movement_label}:} & 3 \\
        \textbf{\trans{strength_label}:} & 3 \\
        \textbf{\trans{dexterity_label}:} & 1 \\
        \textbf{\trans{mind_label}:} & 2 \\
        \textbf{\trans{resistances_label}:} & ['Feuer', 'Gift'] \\
    \end{tabular}

    \textbf{Griff} (6)\\
    Durchschlag 0\\
    \textbf{Fluch} (4)\\
    Durchschlag 3, Gift 2\\


    \subsection*{Vogelscheuche}

    Zur Erntezeit, wenn der Tod in die dämmernde Welt zurückkehrt und die Sommerblüten ihre verwelkten Köpfe neigen, ragen unheimliche Vogelscheuchen in stiller Wache über leeren Feldern auf. Mit unsterblicher Geduld halten diese stoischen Wächter ihre Posten trotz Wind, Sturm und Flut, gebunden an die Befehle ihres Herrn, begierig darauf, Beute mit ihrem Sackleinengesicht zu erschrecken und Opfer mit ihren messerscharfen Klauen zu zerreißen.

    \begin{tabular}{@{}ll@{}}
        \textbf{\trans{wounds_label}:} & 8 \\
        \textbf{\trans{movement_label}:} & 0 \\
        \textbf{\trans{strength_label}:} & 2 \\
        \textbf{\trans{dexterity_label}:} & 0 \\
        \textbf{\trans{mind_label}:} & 4 \\
        \textbf{\trans{resistances_label}:} & ['Feuer'] \\
    \end{tabular}

    \textbf{Klaue} (6)\\
    Blutend 1\\


    \subsection*{Fleischgolem}

    Ein Fleischgolem ist eine grausige Ansammlung humanoider Körperteile, die zu einem muskulösen Ungetüm mit ungeheurer Kraft zusammengenäht und zusammengeschraubt wurden. Mächtige Verzauberungen schützen ihn und wehren Zauber sowie alle Waffen bis auf die stärksten ab.

    \begin{tabular}{@{}ll@{}}
        \textbf{\trans{wounds_label}:} & 10 \\
        \textbf{\trans{movement_label}:} & 4 \\
        \textbf{\trans{strength_label}:} & 2 \\
        \textbf{\trans{dexterity_label}:} & 2 \\
        \textbf{\trans{mind_label}:} & 2 \\
        \textbf{\trans{resistances_label}:} & ['Feuer'] \\
    \end{tabular}

    \textbf{Klaue} (6)\\
    Blutend 1, Vergifted 1\\


    \subsection*{Skelett}

    Ein wandelndes Skelett, von dunkler Magie belebt.

    \begin{tabular}{@{}ll@{}}
        \textbf{\trans{wounds_label}:} & 4 \\
        \textbf{\trans{movement_label}:} & 3 \\
        \textbf{\trans{strength_label}:} & 2 \\
        \textbf{\trans{dexterity_label}:} & 2 \\
        \textbf{\trans{mind_label}:} & 2 \\
    \end{tabular}

    \textbf{Knochengriff} (8)\\
    Vergiftet 1\\
    \textbf{Kalter Atem} (6)\\
    Geschockt 2, 5 Meter Reichweite\\


    \section{Tier}



    \subsection*{Spinne}

    Ein geduldiger und methodischer Jäger, der für Regungslosigkeit, Präzision und das Einfangen steht. Die giftige Spinne wartet verborgen in Netzen oder dunklen Winkeln und schlägt nur zu, wenn ihre Beute gefangen oder wehrlos ist. Ihr Gift schwächt oder lähmt, sodass ein Entkommen kaum möglich ist. Anders als verfolgenden Jägern gehört ihr die Umgebung selbst – sie macht den Raum zur Falle. In düsteren Auslegungen wirkt sie berechnend und fremdartig, eine lautlose Architektin des Todes, deren Präsenz oft erst bemerkt wird, wenn es zu spät ist.

    \begin{tabular}{@{}ll@{}}
        \textbf{\trans{wounds_label}:} & 2 \\
        \textbf{\trans{movement_label}:} & 2 \\
        \textbf{\trans{strength_label}:} & 1 \\
        \textbf{\trans{dexterity_label}:} & 3 \\
        \textbf{\trans{mind_label}:} & 1 \\
    \end{tabular}

    \textbf{Biss} (4)\\
    Vergiftet 1, Durchschlag 1\\


    \subsection*{Mammut}

    Mammut bezeichnet eine Gattung von [[elefant|Elefanten]], die in den kälteren Regionen Tirakans weit verbreitet sind. Die Scholaren unter dem jungen leicht verrückten Zoologen Bernhard von [[yavon|Yavon]] nehmen an das sie sich aus einer Elefantengruppe über mehrere Zwischenformen entwickelten, die sich zunehmend auf Grasnahrung spezialisierten und an die Kälte anpassten.

    \begin{tabular}{@{}ll@{}}
        \textbf{\trans{wounds_label}:} & 18 \\
        \textbf{\trans{movement_label}:} & 6 \\
        \textbf{\trans{strength_label}:} & 6 \\
        \textbf{\trans{dexterity_label}:} & 2 \\
        \textbf{\trans{mind_label}:} & 2 \\
    \end{tabular}

    \textbf{Trampeln} (12)\\
    Durchschlag 0\\
    \textbf{Ansturm} (8)\\
    Durchschlag 0\\


    \subsection*{Riesenkraken}

    Aus der Tiefe entfaltet sich etwas Großes und Geduldiges. Der Riesenkalmar ist ein Meister der Tarnung, verschmilzt mit Fels und Riff, bis er sich entscheidet zu handeln. Seine Intelligenz wirkt beunruhigend – prüfend, tastend, lernend. Wenn er zuschlägt, dann mit plötzlicher Reichweite, umschlingt seine Beute mit kraftvollen Armen und zieht sie in einen erdrückenden, unausweichlichen Griff. Das Wasser selbst wird zur Waffe, verschleiert Sicht und Orientierung. In düsteren Auslegungen wirkt er uralt und wachsam, ein lauernder Verstand in der Tiefe, der auf den richtigen Moment wartet.

    \begin{tabular}{@{}ll@{}}
        \textbf{\trans{wounds_label}:} & 10 \\
        \textbf{\trans{movement_label}:} & 4 \\
        \textbf{\trans{strength_label}:} & 3 \\
        \textbf{\trans{dexterity_label}:} & 3 \\
        \textbf{\trans{mind_label}:} & 2 \\
    \end{tabular}

    \textbf{Würgegriff} (10)\\
    Das Opfer ist gegriffen und erleidet jede Kampfrunde 2 Wunden. Es kann sich in seiner Kampfrunde mit einer Kraft Probe befreien.\\


    \subsection*{Bär}

    Eine gewaltige und widerstandsfähige Naturgewalt, die für rohe Stärke, Ausdauer und territoriale Wut steht. Anders als schnelle Jäger verlässt sich der Bär nicht auf Geschwindigkeit oder Heimlichkeit – einmal aufgebracht, überwältigt er seine Gegner durch schiere Kraft und unerbittliche Härte. Meist wirkt er gleichgültig, doch bei Bedrohung oder zur Verteidigung seines Reviers wird er zu einer verheerenden Macht. In düsteren Auslegungen erscheint er nahezu unaufhaltsam – die Verkörperung urtümlicher Raserei.

    \begin{tabular}{@{}ll@{}}
        \textbf{\trans{wounds_label}:} & 8 \\
        \textbf{\trans{movement_label}:} & 3 \\
        \textbf{\trans{strength_label}:} & 3 \\
        \textbf{\trans{dexterity_label}:} & 2 \\
        \textbf{\trans{mind_label}:} & 2 \\
    \end{tabular}

    \textbf{Biss} (8)\\
    Blutend 2\\
    \textbf{Krallen} (8)\\
    Blutend 1\\


    \subsection*{Krokodil}

    Ein geduldiger, urtümlicher Lauerjäger, der für Regungslosigkeit, Präzision und plötzliche, überwältigende Gewalt steht. Halb im Wasser verborgen, nahezu unsichtbar, schlägt er in einem Augenblick zu, wenn Beute dem Ufer zu nahe kommt. Sein Griff ist unerbittlich – einmal gepackt, gibt es kaum ein Entkommen. Krokodile verfolgen ihre Beute nicht weit; sie beherrschen die Grenze zwischen Land und Wasser und verwandeln sicheren Boden in eine tödliche Falle. In düsteren Auslegungen wirken sie uralt und berechnend, als hätte das Wasser selbst Zähne.

    \begin{tabular}{@{}ll@{}}
        \textbf{\trans{wounds_label}:} & 6 \\
        \textbf{\trans{movement_label}:} & 4 \\
        \textbf{\trans{strength_label}:} & 2 \\
        \textbf{\trans{dexterity_label}:} & 2 \\
        \textbf{\trans{mind_label}:} & 2 \\
    \end{tabular}

    \textbf{Biss} (10)\\
    Das Opfer kann sich nur mit einer Kraft Probe aus dem Biss befreien. Das Krokodil zieht gerne Opfer unter Wasser, wenn dies möglich ist.\\


    \subsection*{Hai}

    Unter der Oberfläche bewegt sich der Hai mit lautloser Sicherheit und nimmt wahr, was unsichtbar bleibt. Er stürzt sich nicht blind auf seine Beute – er kreist, prüft und nähert sich mit kalkulierter Entschlossenheit. Ein plötzlicher Angriff beendet die Begegnung in einem Augenblick aus Bewegung und Blut, bevor er wieder in der Tiefe verschwindet. In seinem Element gibt es keinen festen Halt und kaum ein Entkommen. In düsteren Auslegungen wirkt er wie eine unsichtbare Macht der Tiefe, angezogen von Schwäche und Unruhe – als hätte sich das Wasser selbst gegen seine Beute gewandt.

    \begin{tabular}{@{}ll@{}}
        \textbf{\trans{wounds_label}:} & 8 \\
        \textbf{\trans{movement_label}:} & 6 \\
        \textbf{\trans{strength_label}:} & 2 \\
        \textbf{\trans{dexterity_label}:} & 3 \\
        \textbf{\trans{mind_label}:} & 1 \\
    \end{tabular}

    \textbf{Biss} (10)\\
    Durchschlag 1\\


    \subsection*{Hyene}

    Ein listiger Aasfresser und Gelegenheitsjäger, der dort überlebt, wo andere scheitern. Die Hyäne lebt von Anpassungsfähigkeit, Ausdauer und rücksichtsloser Zweckmäßigkeit und meidet oft den direkten Angriff, bis sich eine Schwäche zeigt. Sie prüfen ihre Beute unermüdlich, zermürben sie und nutzen jede Gelegenheit aus. Ihre unheimlichen Laute und ihr verstörendes Auftreten verleihen ihnen einen beinahe widernatürlichen Ruf. In düsteren Auslegungen wirken sie spöttisch und grausam – Kreaturen, die sich am Verfall und am Untergang anderer laben.

    \begin{tabular}{@{}ll@{}}
        \textbf{\trans{wounds_label}:} & 4 \\
        \textbf{\trans{movement_label}:} & 4 \\
        \textbf{\trans{strength_label}:} & 1 \\
        \textbf{\trans{dexterity_label}:} & 3 \\
        \textbf{\trans{mind_label}:} & 2 \\
    \end{tabular}

    \textbf{Biss} (8)\\
    Blutend 1\\


    \subsection*{Löwe}

    Ein dominanter Jäger, der für Stärke, Autorität und tödliche Präzision steht. Anders als ausdauernde Rudeljäger verlässt er sich auf kurze, überwältigende Angriffe, oft aus dem Hinterhalt, oder behauptet sein Revier allein durch seine Präsenz. Ein Löwe jagt nicht endlos – er entscheidet, wann die Jagd beginnt und endet. In düsteren Auslegungen wirkt er königlich und furchteinflößend, ein Sinnbild von Macht, das keinen Widerspruch duldet.

    \begin{tabular}{@{}ll@{}}
        \textbf{\trans{wounds_label}:} & 6 \\
        \textbf{\trans{movement_label}:} & 4 \\
        \textbf{\trans{strength_label}:} & 3 \\
        \textbf{\trans{dexterity_label}:} & 2 \\
        \textbf{\trans{mind_label}:} & 2 \\
    \end{tabular}

    \textbf{Biss} (8)\\
    Durchschlag 1, Blutend 1\\
    \textbf{Krallen} (6)\\
    Piercing 1, Bleeding 2\\


    \subsection*{Geier}

    Auf aufsteigenden Luftströmen kreist der Geier geduldig und beobachtet jede Regung unter sich. Er jagt nicht im klassischen Sinne, sondern wartet auf das Unvermeidliche, angezogen von Tod, Verfall und leichter Beute. Seine Präsenz wirkt wie ein stummes Omen und erscheint oft schon, bevor das Ende vollständig eingetreten ist. Am Boden bewegt er sich unbeholfen, doch mit unerbittlicher Effizienz beim Reißen und Fressen. In düsteren Auslegungen wirkt er wie ein Vorbote des Schicksals, der überall dort erscheint, wo Leben schwindet – als würde ihn das Sterben selbst rufen.

    \begin{tabular}{@{}ll@{}}
        \textbf{\trans{wounds_label}:} & 2 \\
        \textbf{\trans{movement_label}:} & 6 \\
        \textbf{\trans{strength_label}:} & 1 \\
        \textbf{\trans{dexterity_label}:} & 3 \\
        \textbf{\trans{mind_label}:} & 2 \\
    \end{tabular}

    \textbf{Schnabel} (7)\\
    Durchschlag 1\\


    \subsection*{Robbe}

    An kalten Küsten und treibendem Eis bewegt sich die Robbe mühelos zwischen Wasser und Land. Im Meer ist sie schnell und schwer fassbar, taucht und wendet sich mit fließender Präzision; an Land wirkt sie langsamer, bleibt jedoch wachsam und vorsichtig. Robben sind keine von Natur aus aggressiven Tiere, doch in die Enge getrieben oder beim Schutz ihres Nachwuchses reagieren sie mit plötzlicher Heftigkeit – beißen und schlagen mit überraschender Kraft um sich. Ihr Auftreten weist oft auf reiche Jagdgründe hin, aber auch auf verborgene Gefahren unter der Oberfläche. In düsteren Auslegungen wirken sie wachsam und wissend, ihre Augen knapp über der Wasserlinie, bevor sie lautlos verschwinden.

    \begin{tabular}{@{}ll@{}}
        \textbf{\trans{wounds_label}:} & 5 \\
        \textbf{\trans{movement_label}:} & 5 \\
        \textbf{\trans{strength_label}:} & 1 \\
        \textbf{\trans{dexterity_label}:} & 2 \\
        \textbf{\trans{mind_label}:} & 2 \\
    \end{tabular}

    \textbf{Biss} (5)\\
    Durchschlag 0\\


    \subsection*{Nashorn}

    Ein massiger und reizbarer Gigant, der für rohe Wucht, unbeugsame Zähigkeit und plötzliche Aggression steht. Sein schwaches Sehvermögen wird durch feines Gehör und Geruchssinn ausgeglichen, sodass es auf jede wahrgenommene Bedrohung sofort reagiert. Hat es sich einmal entschieden, stürmt es mit unaufhaltsamer Kraft voran und spießt oder zertrampelt alles auf seinem Weg. Ein Nashorn weicht nicht aus und zögert nicht – es durchbricht. In düsteren Auslegungen wirkt es wie ein lebender Rammbock, getrieben von blinder Wucht und unerschütterlichem Instinkt.

    \begin{tabular}{@{}ll@{}}
        \textbf{\trans{wounds_label}:} & 10 \\
        \textbf{\trans{movement_label}:} & 4 \\
        \textbf{\trans{strength_label}:} & 3 \\
        \textbf{\trans{dexterity_label}:} & 3 \\
        \textbf{\trans{mind_label}:} & 1 \\
    \end{tabular}

    \textbf{Ansturm} (10)\\
    Durchschlag 2\\


    \subsection*{Weiterschnatterer}

    Im 7. Jahrhundert wird man durch die [[gnome|Gnome]] auf diesen seltsamen Vogel, der auf kleinen Inseln im Felsenmeer lebt auch militärisch aufmerksam. Gelb und grün ist sein Federvieh und seltsam krumm sein kurzer Schnabel. Das besondere an ihnen ist, das wenn mehrere von ihnen nebeinander sind und man dem einen ein Wort zuflüstert, es automatisch alle anderen Vögel mit dem Angesprochenen wiederholen.

    \begin{tabular}{@{}ll@{}}
        \textbf{\trans{wounds_label}:} & 2 \\
        \textbf{\trans{movement_label}:} & 11 \\
        \textbf{\trans{strength_label}:} & 1 \\
        \textbf{\trans{dexterity_label}:} & 4 \\
        \textbf{\trans{mind_label}:} & 7 \\
    \end{tabular}



    \subsection*{Dragsberger Scheunenteufel}

    Die Dragsberger Scheunenteufel sind wohl größte Hauskatzenart Tirakans und stammen aus dem gleichnamigen Herzogtum Dragsberg in [[asgoran|Asgoran]]. Ein Dragsberger Scheunenteufel wird 1 bis 1,40 Schritt lang. Sie haben ein 3-4 Finger langes Fell von unterschiedlichsten Farben, 2 und 3-farbige Tiere gehören zur Regel, und einen langen buschigen Schwanz.

    \begin{tabular}{@{}ll@{}}
        \textbf{\trans{wounds_label}:} & 3 \\
        \textbf{\trans{movement_label}:} & 6 \\
        \textbf{\trans{strength_label}:} & 2 \\
        \textbf{\trans{dexterity_label}:} & 4 \\
        \textbf{\trans{mind_label}:} & 3 \\
        \textbf{\trans{resistances_label}:} & ['Feuer', 'Schock'] \\
    \end{tabular}

    \textbf{Kratzen} (4)\\
    Die Krallen des Scheunenteufels sind sehr unangenehm, bekommt man sie durch das Gesicht gezogen.\\
    \textbf{Beissen} (5)\\
    Durchschlag 0\\


    \subsection*{Hochseeschlange}

    Eine widerliche Schlangenart, die vor allem an der Ostküste, um die [[echsen|Echseninsel]] herum auftritt und dort beinahe die komplette Wasserlandschaft übernommen hat. Sie ist etwa 3-6 Schritt lang und schimmert leicht grünlich. Sie hat Kiemen an ihrem sonst schlangenartigen Kopf, das Maul zieren sich ineinander verhakende Fangzähne.

    \begin{tabular}{@{}ll@{}}
        \textbf{\trans{wounds_label}:} & 15 \\
        \textbf{\trans{movement_label}:} & 8 \\
        \textbf{\trans{strength_label}:} & 5 \\
        \textbf{\trans{dexterity_label}:} & 4 \\
        \textbf{\trans{mind_label}:} & 2 \\
    \end{tabular}

    \textbf{Biss} (10)\\
    Durchschlag 1, Blutend 1\\


    \subsection*{Brauner Tiger}

    Weit entfernt von seinem Vetter, dem gelben Tiger, lebt der braune Tiger im höchsten Norden Tirakans. Ihn unterscheidet vorallem das insgesamt dunkle Fell von allen anderen Tigern. Weiterhin besitzt er gewaltige Eckzähne, welche die größten im Reich der Raubkatzen sind.

    \begin{tabular}{@{}ll@{}}
        \textbf{\trans{wounds_label}:} & 8 \\
        \textbf{\trans{movement_label}:} & 7 \\
        \textbf{\trans{strength_label}:} & 2 \\
        \textbf{\trans{dexterity_label}:} & 4 \\
        \textbf{\trans{mind_label}:} & 4 \\
        \textbf{\trans{resistances_label}:} & ['Kälte'] \\
    \end{tabular}

    \textbf{Krallen} (10)\\
    Durchschlag 1\\
    \textbf{Beissen} (10)\\
    Durchschlag 1, Blutend 1\\


    \subsection*{Ziege}

    Dieses Tier, dass die Scholaren in die Familie der sogenannten Huftiere eingeordnet haben (erste Erwähnung im alten Codex Bestiarius im [[viertes-zeitalter-das-zeitalter-der-menschen|ersten Jahrtausend vor dem neuen Zeitalter]]) ist verwandt mit dem im Norden bekannteren Schaf. Forschungen neuester Zeit zeigen daß sich die Ziege wohl im Süden Tirakans aus eben jenen Schafen zu entwickelt haben scheint (eine These die besonders bei den religiösen Scholaren hart umstritten ist).

    \begin{tabular}{@{}ll@{}}
        \textbf{\trans{wounds_label}:} & 4 \\
        \textbf{\trans{movement_label}:} & 4 \\
        \textbf{\trans{strength_label}:} & 3 \\
        \textbf{\trans{dexterity_label}:} & 2 \\
        \textbf{\trans{mind_label}:} & 2 \\
    \end{tabular}

    \textbf{Treten} (6)\\
    Durchschlag 0\\
    \textbf{Biss} (8)\\
    Durchschlag 0\\


    \subsection*{Sumpfturtoll}

    So werden 2 Schritt große, gelbäugige Kröten genannt. Diese mystischen Tiere sind genauso selten wie gefährlich. In ihrem unersättlichem Hunger stürzen sie sich auf jegliche Lebensform, die ihnen über den Weg läuft und nicht doppelt so massig wie sie selbst ist. Um sein Opfer zu töten bespeiht der Sumpfturtoll es mit einer dickflüssigen, stark ätzenden Säure, welche es selbst vermag das Eisen der Rüstungen zu zerstören.

    \begin{tabular}{@{}ll@{}}
        \textbf{\trans{wounds_label}:} & 4 \\
        \textbf{\trans{movement_label}:} & 2 \\
        \textbf{\trans{strength_label}:} & 1 \\
        \textbf{\trans{dexterity_label}:} & 3 \\
        \textbf{\trans{mind_label}:} & 5 \\
        \textbf{\trans{resistances_label}:} & ['Gift'] \\
    \end{tabular}

    \textbf{Spucken} (8)\\
    Die Kröte spuckt eine gelbliche ätzende Flüssigkeit in einem Winkel von 90 Grad 5 Schritt weit. Die Flüssigkeit verursacht eine Wunde pro Kampfrunde bis sie abgewaschen wurde. Zudem verringert sie den Zustand aller betroffenen Rüstungsteile um 1 Schutz.\\


    \subsection*{Snake}

    Ein lautloser und schwer fassbarer Jäger, der für Geduld, Präzision und verborgene Tödlichkeit steht. Die Giftschlange verlässt sich auf Tarnung und regungslose Ruhe, um blitzschnell zuzuschlagen, sobald sich Beute nähert. Ihr Gift übernimmt den eigentlichen Tod – es schwächt, lähmt oder zerstört das Opfer langsam von innen. Schlangen meiden offene Konfrontationen und setzen auf Hinterhalt und Rückzug, doch ein einziger Fehler kann tödlich sein. In düsteren Auslegungen wirken sie kalt und berechnend, als sei jede Bewegung geplant und jeder Biss unausweichlich.

    \begin{tabular}{@{}ll@{}}
        \textbf{\trans{wounds_label}:} & 2 \\
        \textbf{\trans{movement_label}:} & 5 \\
        \textbf{\trans{strength_label}:} & 1 \\
        \textbf{\trans{dexterity_label}:} & 3 \\
        \textbf{\trans{mind_label}:} & 1 \\
    \end{tabular}

    \textbf{Biss} (8)\\
    Vergiftet 1\\


    \subsection*{Silbermondfisch}

    Diese Fischart ist extrem selten, einzelgängerisch und nur auf hoher See anzutreffen. Der Leib dieses Fisches ist seltsam scheibenartig gestaltet, mit einem Durchmesser von bis zu zwei Schritt, die Flossen sind rund. Seine Schuppen haben eine hellblaue Farbe, sind aber fast durchsichtig.

    \begin{tabular}{@{}ll@{}}
        \textbf{\trans{wounds_label}:} & 4 \\
        \textbf{\trans{movement_label}:} & 7 \\
        \textbf{\trans{strength_label}:} & 1 \\
        \textbf{\trans{dexterity_label}:} & 5 \\
        \textbf{\trans{mind_label}:} & 1 \\
    \end{tabular}



    \subsection*{Ratte}

    In Dunkelheit und verborgenen Winkeln gedeiht die Ratte dort, wo andere scheitern. Sie zwängt sich durch kleinste Spalten, nistet sich in Wänden ein und tritt oft in solcher Zahl auf, dass aus einer Plage eine Bedrohung wird. Allein ist sie scheu und leicht zu vertreiben, doch im Schwarm wird sie kühn – sie nagt, beißt und überwältigt durch schiere Hartnäckigkeit. Ratten verbreiten Schmutz, vernichten Vorräte und bringen unsichtbare Gefahren mit sich. In düsteren Auslegungen wirken sie wie eine schleichende Verseuchung, ein lebendiges Zeichen des Verfalls, das stärker wird, je länger man es ignoriert.

    \begin{tabular}{@{}ll@{}}
        \textbf{\trans{wounds_label}:} & 2 \\
        \textbf{\trans{movement_label}:} & 5 \\
        \textbf{\trans{strength_label}:} & 1 \\
        \textbf{\trans{dexterity_label}:} & 3 \\
        \textbf{\trans{mind_label}:} & 1 \\
    \end{tabular}

    \textbf{Biss} (5)\\
    Blutend 1\\


    \subsection*{Skorpion}

    Unter Sand, Stein oder im Schatten verborgen, verharrt der Skorpion in unheimlicher Regungslosigkeit. Wird er gestört, reagiert er augenblicklich – sein Stachel trifft mit präziser, geübter Geschwindigkeit. Seine Stärke liegt nicht in roher Gewalt, sondern im Gift, das schwächt, lähmt oder den Kampf sofort beendet. Er verfolgt nicht – er bestraft das Eindringen. In düsteren Auslegungen wirkt er uralt und gleichgültig, perfekt angepasst an eine Welt, in der Zögern den Tod bedeutet.

    \begin{tabular}{@{}ll@{}}
        \textbf{\trans{wounds_label}:} & 3 \\
        \textbf{\trans{movement_label}:} & 5 \\
        \textbf{\trans{strength_label}:} & 1 \\
        \textbf{\trans{dexterity_label}:} & 3 \\
        \textbf{\trans{mind_label}:} & 1 \\
    \end{tabular}

    \textbf{Biss} (6)\\
    Vergiftet 1\\
    \textbf{Stachel} (8)\\
    Vergiftet 2\\


    \subsection*{Hadidim-Echse}

    Die Hadidim-Echsen sind große Flugechsen mit bis zu 2 Schritt Spannweite. Sie ernähren sich vor allem von Fischen, sind aber auch dem einen oder anderen Meeresvogel nicht abgeneigt. Hadidim-Echsen kommen nur am schmalen Küstenstreifen und an den Seen und Oasen [[al-bah-jira|Al Bah JiRas]] vor. [[menschen|Menschen]] werden sie nur sehr selten gefährlich,es sei denn jemand ist so töricht und versucht einer Hadidim-Mutter die Eier aus dem Nest zu stehlen.

    \begin{tabular}{@{}ll@{}}
        \textbf{\trans{wounds_label}:} & 14 \\
        \textbf{\trans{movement_label}:} & 12 \\
        \textbf{\trans{strength_label}:} & 3 \\
        \textbf{\trans{dexterity_label}:} & 3 \\
        \textbf{\trans{mind_label}:} & 2 \\
    \end{tabular}

    \textbf{Biss} (8)\\
    Durchschlag 1\\


    \subsection*{Frosch}

    Leuchtende Farben im Unterholz warnen vor Gefahr, lange bevor sich der Frosch selbst bewegt. Klein und leicht zu übersehen, verlässt er sich nicht auf Stärke, sondern auf hochwirksame Gifte, die selbst eine unachtsame Berührung gefährlich machen. Er jagt nicht durch Verfolgung, sondern durch Nähe – seine bloße Anwesenheit kann bereits schaden. In dichten Dschungeln oder feuchten Umgebungen wird er Teil der Landschaft, eine verborgene Gefahr zwischen Blättern und Wasser. In düsteren Auslegungen wirkt er beinahe unnatürlich, seine grelle Erscheinung wie eine bewusste Warnung, dass unter der Oberfläche etwas nicht stimmt.

    \begin{tabular}{@{}ll@{}}
        \textbf{\trans{wounds_label}:} & 2 \\
        \textbf{\trans{movement_label}:} & 4 \\
        \textbf{\trans{strength_label}:} & 1 \\
        \textbf{\trans{dexterity_label}:} & 1 \\
        \textbf{\trans{mind_label}:} & 1 \\
    \end{tabular}

    \textbf{Zunge} (3)\\
    Vergiftet 1\\


    \subsection*{Tiger}

    Ein einzelgängerischer und schwer fassbarer Prädator, der für Heimlichkeit, Geduld und plötzliche, explosive Gewalt steht. Anders als Rudeljäger verlässt er sich auf den perfekten Moment: ungesehen anschleichen, lautlos nähern und die Jagd mit einem einzigen, entscheidenden Angriff beenden. Tiger sind territorial und unberechenbar, verschwinden in ihrer Umgebung ebenso mühelos, wie sie daraus hervortreten. In düsteren Auslegungen wirken sie beinahe geisterhaft – lautlos, unausweichlich und tödlich präzise.

    \begin{tabular}{@{}ll@{}}
        \textbf{\trans{wounds_label}:} & 6 \\
        \textbf{\trans{movement_label}:} & 4 \\
        \textbf{\trans{strength_label}:} & 2 \\
        \textbf{\trans{dexterity_label}:} & 3 \\
        \textbf{\trans{mind_label}:} & 2 \\
    \end{tabular}

    \textbf{Biss} (10)\\
    Durchschlag 1, Blutend 1\\
    \textbf{Krallen} (10)\\
    Durchschlag 1\\


    \subsection*{Qualle}

    Lautlos und beinahe anmutig treibt die Qualle durch das Wasser und wirkt zunächst harmlos – bis es zur Berührung kommt. Ihr durchscheinender Körper verbirgt lange Tentakel, die mit starkem Gift versehen sind und augenblicklich reagieren. Sie jagt nicht im klassischen Sinne, sondern macht ihre Umgebung selbst zur Gefahr, in der jede Bewegung riskant wird. In Gruppen bilden sie fast unsichtbare Barrieren, die oft erst bemerkt werden, wenn es zu spät ist. In düsteren Auslegungen wirken sie wie lebende Fallen des Meeres – still, gleichgültig und unausweichlich.

    \begin{tabular}{@{}ll@{}}
        \textbf{\trans{wounds_label}:} & 2 \\
        \textbf{\trans{movement_label}:} & 4 \\
        \textbf{\trans{strength_label}:} & 1 \\
        \textbf{\trans{dexterity_label}:} & 2 \\
        \textbf{\trans{mind_label}:} & 1 \\
    \end{tabular}

    \textbf{Stachel} (4)\\
    Vergiftet 1\\


    \subsection*{Sumpfkrähe}

    Dies ist eine außergewöhnlich großwüchsige Variante der weithin bekannten gemeinen Krähe. Ihre Heimat sind die Sümpfe im Osten [[yadosien|Yadosiens]]. Wanderer, die an diesem Ort unterwegs sind sollten auf ihr Gepäck achtgeben, denn es kann leicht passieren, dass in einem unachtsamen Moment eine Sumpfkrähe vom Himmel herabgleitet, es an sich reisst und in ihr Versteck trägt.

    \begin{tabular}{@{}ll@{}}
        \textbf{\trans{wounds_label}:} & 2 \\
        \textbf{\trans{movement_label}:} & 7 \\
        \textbf{\trans{strength_label}:} & 1 \\
        \textbf{\trans{dexterity_label}:} & 4 \\
        \textbf{\trans{mind_label}:} & 4 \\
    \end{tabular}

    \textbf{Stehlen} (8)\\
    Gelingt der Krähe der Wurf, so kann der Bestohlene einen Wahrnehmungswurf machen, dessen Mindestwurf um die Anzahl der Erfolge der Krähe angehoben ist. Bei Gelingen des Wahrnehmungswurfs wird die Krähe auf frischer Tat ertappt, ansonsten gelingt ihr der Diebstahl unbemerkt.\\
    \textbf{Picken} (6)\\
    Blutend 1\\


    \subsection*{Baumreisser}

    Ein ausgewachsenes Tier misst etwa drei Schritt Länge und wiegt etwa fünf Zentner. Seine kräftigen Beine machen es zu einem hervoragenden Kletterer und schnellen Läufer. Der Baumreisser ernährt sich hauptsächlich von Kleintieren und Aas, schreckt in schlechten Zeiten aber auch nicht vor Angriffen auf andere Kreaturen zurück. Grundsätzlich ist jedes Wesen, das kleiner ist als es selbst, eine mögliche Erweiterung des Speiseplans.

    \begin{tabular}{@{}ll@{}}
        \textbf{\trans{wounds_label}:} & 10 \\
        \textbf{\trans{movement_label}:} & 6 \\
        \textbf{\trans{strength_label}:} & 4 \\
        \textbf{\trans{dexterity_label}:} & 2 \\
        \textbf{\trans{mind_label}:} & 2 \\
    \end{tabular}

    \textbf{Krallen} (5)\\
    Durchschlag 1\\
    \textbf{Würgen} (10)\\
    Mit seinen starken Armen ist der Baumreisser in der Lage sein Opfer zu würgen und sich zudem noch zu bewegen. Das Würgen hat Durchschlag 0. Es bedarf einer schweren Kraftprobe, um sich aus dem Griff zu befreien.\\


    \subsection*{Höllentaipan}

    Der Höllentaipan ist eine schlanke Schlange von bis zu 3 Schritt Länge, von dunkelbrauner, seltener olivgrüner, Farbe. Er kommt in den südlichen Gebieten von [[yadosien|Yadosien]], über [[al-bah-jira|Al Bah JiRa]] bis hin zu den Landen der [[ogrut|O'Grut]] und den [[ancatir|Ancatir]] vor.

    \begin{tabular}{@{}ll@{}}
        \textbf{\trans{wounds_label}:} & 3 \\
        \textbf{\trans{movement_label}:} & 2 \\
        \textbf{\trans{strength_label}:} & 2 \\
        \textbf{\trans{dexterity_label}:} & 4 \\
        \textbf{\trans{mind_label}:} & 4 \\
        \textbf{\trans{resistances_label}:} & ['Gift'] \\
    \end{tabular}

    \textbf{Biss} (8)\\
    Durchschlag 4, Vergiften 2\\


    \subsection*{Große Höhlenspinne}

    Die große Höhlenspinne kommt in allen Gebieten des zentralen Tirakans vor. Sie erreicht eine Höhe von 3 Schritt, und sie ist trotz ihres Gewichtes äußerst geschickt dabei, sich durch Gänge zu bewegen.

    \begin{tabular}{@{}ll@{}}
        \textbf{\trans{wounds_label}:} & 10 \\
        \textbf{\trans{movement_label}:} & 4 \\
        \textbf{\trans{strength_label}:} & 2 \\
        \textbf{\trans{dexterity_label}:} & 4 \\
        \textbf{\trans{mind_label}:} & 3 \\
        \textbf{\trans{resistances_label}:} & ['Feuer', 'Gift'] \\
    \end{tabular}

    \textbf{Biss} (8)\\
    Vergiftet 1\\
    \textbf{Netz} (8)\\
    Geschockt 1\\


    \subsection*{Schaf}

    Schafe sind meist wollige, manchmal bis knapp über ein Schritt hohe Nutztiere. Ihre Farbe schwankt von weiß über grau bis schwarz. Neben den Wollrassen gibt es auch Rassen mit weniger Wolle welche in erster Linie zur Fleischproduktion gezüchtet werden.

    \begin{tabular}{@{}ll@{}}
        \textbf{\trans{wounds_label}:} & 5 \\
        \textbf{\trans{movement_label}:} & 4 \\
        \textbf{\trans{strength_label}:} & 2 \\
        \textbf{\trans{dexterity_label}:} & 1 \\
        \textbf{\trans{mind_label}:} & 2 \\
    \end{tabular}

    \textbf{Blöken} (10)\\
    Das Blöken einer ganzen Herde Schafe kann für einen Charakter eine ziemliche Verwirrung bedeuten.\\
    \textbf{Treten} (6)\\
    Durchschlag 0\\
    \textbf{Biss} (8)\\
    Durchschlag 0\\


    \subsection*{Sandviper}

    Diese garstigen Schlangen bewohnen die große Wüste [[al-bah-jira|Al Bah JiRa]]. Zumeist liegen sie verborgen im Sand, wo sie vorbeigehenden Tieren oder auch Menschen auflauern. Wenn dann tatsächlich ein Unglücklicher in Reichweite einer lauernden Sandviper geraten sollte, so wird er erleben, wie scheinbar aus dem nichts ein schuppiges, gehörntes Gräuel mit einer Länge von 2 Schritt und 1 Fuß Dicke emporschnellen und sich an ihm festbeißen wird. Das hochwirksame Gift wird ihn innerhalb von vier Stunden töten, sollte es ihm bis dahin nicht gelingen ein Gegengift zu bekommen.

    \begin{tabular}{@{}ll@{}}
        \textbf{\trans{wounds_label}:} & 3 \\
        \textbf{\trans{movement_label}:} & 2 \\
        \textbf{\trans{strength_label}:} & 1 \\
        \textbf{\trans{dexterity_label}:} & 5 \\
        \textbf{\trans{mind_label}:} & 4 \\
    \end{tabular}

    \textbf{Biss} (8)\\
    Vergiften 2\\


    \subsection*{Wolf}

    Überall in den gemäßigten Regionen Tirakans gibt es Wölfe. Diese Rudeltiere sind in der Regel zurückhaltend, und halten sich von Menschen fern. Wenn sie in die Enge getrieben oder bedroht werden, werden sie aggresiv und können auch Menschen angreifen.

    \begin{tabular}{@{}ll@{}}
        \textbf{\trans{wounds_label}:} & 5 \\
        \textbf{\trans{movement_label}:} & 7 \\
        \textbf{\trans{strength_label}:} & 2 \\
        \textbf{\trans{dexterity_label}:} & 2 \\
        \textbf{\trans{mind_label}:} & 2 \\
    \end{tabular}

    \textbf{Biss} (7)\\
    Bluten 1\\
    \textbf{Ruf nach dem Rudel} (4)\\
    Der Wolf heult laut auf. Für jeden Erfolg dieser Aktion erscheinen in 1W3 Kampfrunden weitere Wölfe in Anzahl der Erfolge des Wurfes, falls weitere Wölfe in der Nähe sind.\\


    \subsection*{Nahrz'gu}

    Die Nahrz'gu sind Pflanzenfresser, die in den nördlichen Ebenen und Wäldern vorkommen. Sie erreichen ein Schulterhöhe von knapp 1,7 Schritt und sind eine kräftige Abart des Auerochsens. Sie haben ein langes struppiges Fell und dicke nach vorn gebogene Hörner.

    \begin{tabular}{@{}ll@{}}
        \textbf{\trans{wounds_label}:} & 10 \\
        \textbf{\trans{movement_label}:} & 7 \\
        \textbf{\trans{strength_label}:} & 4 \\
        \textbf{\trans{dexterity_label}:} & 2 \\
        \textbf{\trans{mind_label}:} & 2 \\
    \end{tabular}

    \textbf{Ansturm} (10)\\
    Schock 1\\


    \subsection*{Schwarmfresser}

    Der Schwarmfresser ist vor allem in den flachen Regionen des großen Binnensees von [[yavon|Yavon]] heimisch. Doch auch in anderen Bereichen des langen Flusses findet man einzelne Exemplare. Ausgewachsen erreichen diese Fische etwa die Länge von einem Schritt. Sie kennzeichnen sich durch eine grün-braun gestreifte Oberseite und eine silbrige Unterseite. Die Heckflosse ist oben leicht abgerundet.

    \begin{tabular}{@{}ll@{}}
        \textbf{\trans{wounds_label}:} & 3 \\
        \textbf{\trans{movement_label}:} & 6 \\
        \textbf{\trans{strength_label}:} & 2 \\
        \textbf{\trans{dexterity_label}:} & 3 \\
        \textbf{\trans{mind_label}:} & 1 \\
    \end{tabular}



    \subsection*{Ochse}

    Ein kräftiges und eigensinniges Tier, das für rohe Gewalt, Ausdauer und territoriale Wut steht. Ob wild oder in Rage versetzt, reagiert es mit plötzlichen, heftigen Angriffen und setzt seine Hörner ein, um Gegner aufzuschlitzen oder fortzuschleudern. Anders als Raubtiere jagt es nicht – es stellt sich entgegen und vertreibt Eindringlinge durch schiere Wucht und Einschüchterung. Ochsen sind zwar kontrollierter, verfügen jedoch über enorme Kraft und können unter Stress oder Überlastung gefährlich werden. In düsteren Auslegungen wird der Bulle zum Sinnbild blinder Raserei und unaufhaltsamer Wucht – eine lebende Kraft, die sich nicht aufhalten lässt.

    \begin{tabular}{@{}ll@{}}
        \textbf{\trans{wounds_label}:} & 8 \\
        \textbf{\trans{movement_label}:} & 4 \\
        \textbf{\trans{strength_label}:} & 2 \\
        \textbf{\trans{dexterity_label}:} & 3 \\
        \textbf{\trans{mind_label}:} & 2 \\
    \end{tabular}

    \textbf{Rammen} (10)\\
    Durchschlag 1\\


    \subsection*{Kinstarchel}

    Kinstarchel werden in [[hadewald|Hadewald]] wild lebende Zwergrinder genannt. Die größten Exemplare dieser Art erreichen selten eine Schulterhöhe von mehr als 2 Ellen. Die kleinen stämmigen Tiere sind silbergrau bis pechschwarz gefärbt und ziehen in Rudeln durch menschenabgelegene Waldstücke. Kinstarchel sind allzeit beliebte Jagdziele. Weil sie immer seltener werden, sind in letzter Zeit Bemühungen gemacht worden, Kinstarchel zu domestizieren. Dies erweist sich jedoch bisher als schwierig, da die gefangenen Exemplare oft nach wenigen Monaten an einer seltsamen Krankheit sterben.

    \begin{tabular}{@{}ll@{}}
        \textbf{\trans{wounds_label}:} & 6 \\
        \textbf{\trans{movement_label}:} & 3 \\
        \textbf{\trans{strength_label}:} & 2 \\
        \textbf{\trans{dexterity_label}:} & 1 \\
        \textbf{\trans{mind_label}:} & 2 \\
    \end{tabular}



    \subsection*{Elefant}

    Ein gewaltiger Riese der Natur, der für enorme Stärke, Widerstandskraft und unaufhaltsame Wucht steht. Meist ruhig und bedacht, wird er bei Bedrohung oder Provokation zu einer verheerenden Macht – er stürmt mit erdrückendem Gewicht voran und fegt alles beiseite, was sich ihm in den Weg stellt. Schon seine Größe verändert das Schlachtfeld und macht das Gelände zu seinem Vorteil. Elefanten sind nicht von Natur aus aggressiv, doch einmal erzürnt oder in Panik, sind sie kaum aufzuhalten. In düsteren Auslegungen wirken sie wie lebende Zerstörungsmaschinen – schwerfällig im Zorn, aber vernichtend in ihrer Wirkung.

    \begin{tabular}{@{}ll@{}}
        \textbf{\trans{wounds_label}:} & 10 \\
        \textbf{\trans{movement_label}:} & 4 \\
        \textbf{\trans{strength_label}:} & 4 \\
        \textbf{\trans{dexterity_label}:} & 2 \\
        \textbf{\trans{mind_label}:} & 2 \\
    \end{tabular}

    \textbf{Würgen} (8)\\
    Der Elefant würgt das Opfer mit seinem Rüssel.\\
    \textbf{Stampfen} (16)\\
    Das ganze Gewicht des Elefanten lastet auf dem Opfer. Der Angriff verursacht Durchschlag 2, kann aber nur ausgeführt werden ,wenn der Elefant bereits eine Runde neben dem Opfer steht.\\


    \subsection*{Adler}

    Hoch über dem Land zieht der Adler seine Kreise und überblickt mit scharfem Blick jede Bewegung unter sich. Er verschwendet keine Kraft – er wartet auf den richtigen Moment, bevor er in einem schnellen, tödlichen Sturzflug angreift. Seine Fänge treffen mit Präzision und Wucht und greifen die Beute, bevor sie reagieren kann. Selbst ohne Angriff beherrscht seine Präsenz den Himmel, ein ständiger Beobachter, dem man sich kaum entziehen kann. In düsteren Auslegungen wirkt er wie ein unerbittliches Auge des Himmels – fern, geduldig und absolut entschlossen, wenn er zuschlägt.

    \begin{tabular}{@{}ll@{}}
        \textbf{\trans{wounds_label}:} & 2 \\
        \textbf{\trans{movement_label}:} & 10 \\
        \textbf{\trans{strength_label}:} & 1 \\
        \textbf{\trans{dexterity_label}:} & 4 \\
        \textbf{\trans{mind_label}:} & 2 \\
    \end{tabular}

    \textbf{Krallen} (7)\\
    Blutend 1\\
    \textbf{Schnabel} (7)\\
    Durchschlag 1\\


    \subsection*{Augura}

    Die Augura sind eulenartige, riesenhafte Vögel. Obwohl sie nicht magisch sind, können sie aufgrund ihrer Größe mit ihren Krallen und ihrem Schnabel beängstigende Wunden reißen. Sie leben in den schneebedeckten Ebenen Bitrheimrs. Augura haben sehr feine Sinne und werden bei der kleinsten Störung ihrer Ruhe aufgescheucht. Sie werden sehr schnell aggressiv.

    \begin{tabular}{@{}ll@{}}
        \textbf{\trans{wounds_label}:} & 4 \\
        \textbf{\trans{movement_label}:} & 7 \\
        \textbf{\trans{strength_label}:} & 2 \\
        \textbf{\trans{dexterity_label}:} & 5 \\
        \textbf{\trans{mind_label}:} & 4 \\
    \end{tabular}

    \textbf{Biss} (6)\\
    Durchschlag 1\\
    \textbf{Krallen} (7)\\
    Durchschlag 1, Blutend 1\\
    \textbf{Ausweichen} (4)\\
    Siehe Ausweichen Spielerregel\\


    \subsection*{Fledermaus}

    Von Dachsparren und Höhlendecken stoßen Fledermäuse in rastloser Bewegung herab, geleitet eher von Klang als von Sicht. In unberechenbaren Schwärmen wirbeln sie durch die Luft, verwirren und beunruhigen, und ihre plötzliche Nähe kann selbst Erfahrene erschrecken. Allein sind sie zerbrechlich, doch in der Masse überwältigen sie – kratzen, beißen und erfüllen die Luft mit chaotischer Bewegung. Ihr Auftreten kündet oft von Verfall, Dunkelheit oder vergessenen Orten. In düsteren Auslegungen wirken sie wie lebendige Schatten, die aus der Finsternis hervorbrechen und ebenso schnell wieder verschwinden.

    \begin{tabular}{@{}ll@{}}
        \textbf{\trans{wounds_label}:} & 2 \\
        \textbf{\trans{movement_label}:} & 5 \\
        \textbf{\trans{strength_label}:} & 1 \\
        \textbf{\trans{dexterity_label}:} & 4 \\
        \textbf{\trans{mind_label}:} & 1 \\
    \end{tabular}

    \textbf{Biss} (2)\\
    Durchschlag 2\\


    \subsection*{Wildschwein}

    Eine brutale und furchtlose Naturgewalt, die für Aggression, Zähigkeit und territoriale Dominanz steht. Wird ein Wildschwein bedroht oder in die Enge getrieben, greift es mit explosiver Geschwindigkeit an und setzt seine Hauer ein, um Gegner aufzuschlitzen und niederzuwalzen. Wildschweine weichen kaum zurück – sie kämpfen kompromisslos weiter, selbst verletzt und im Schmerz. In düsteren Auslegungen wirken sie nahezu unaufhaltsam, getrieben von blinder Wut und dem unerbittlichen Instinkt, Eindringlinge zu vernichten.

    \begin{tabular}{@{}ll@{}}
        \textbf{\trans{wounds_label}:} & 8 \\
        \textbf{\trans{movement_label}:} & 4 \\
        \textbf{\trans{strength_label}:} & 3 \\
        \textbf{\trans{dexterity_label}:} & 1 \\
        \textbf{\trans{mind_label}:} & 2 \\
    \end{tabular}

    \textbf{Rammen} (10)\\
    Durchschlag 1\\


    \subsection*{Dreihörniges Nashorn}

    Angeblich, wenn man diesen neumodischen jungen Scholaren glaubt, welche sich mit eigenartigen Fächern wie Zoologie beschäftigen, ist das dreihörnige Nashorn ein Verwandter des Elefanten und des Mammuts. Doch seht Euch dieses Tier an! Kein Rüssel! Wesentliches Merkmal der Tiere sind die 3 Hörner. Das vordere Horn entwächst dem Nasenbein, das mittlere und das hintere dem Vorderschädel.

    \begin{tabular}{@{}ll@{}}
        \textbf{\trans{wounds_label}:} & 12 \\
        \textbf{\trans{movement_label}:} & 8 \\
        \textbf{\trans{strength_label}:} & 4 \\
        \textbf{\trans{dexterity_label}:} & 2 \\
        \textbf{\trans{mind_label}:} & 2 \\
    \end{tabular}

    \textbf{Ansturm} (10)\\
    Durchschlag 1\\
    \textbf{Trampeln} (10)\\
    Durchschlag 0\\


    \subsection*{Kamel}

    Ein widerstandsfähiges Tier rauer Landschaften, das in eine aggressive und gefährliche Form umschlägt. Kamele sind normalerweise zäh und bedächtig, können jedoch unter Stress oder in ihrem Revier äußerst unberechenbar werden und mit kräftigen Bissen, Tritten und plötzlichen Angriffen reagieren. Ihre Ausdauer erlaubt es ihnen, Gegner zu überdauern und Angriffe länger aufrechtzuerhalten, als man erwarten würde. Einmal erzürnt, lassen sie sich kaum einschüchtern – ihre Sturheit schlägt in unerbittliche Feindseligkeit um. In düsteren Auslegungen wirken sie unheimlich – zäh, nachtragend und beunruhigend hartnäckig, als würde die Wüste selbst sich wehren.

    \begin{tabular}{@{}ll@{}}
        \textbf{\trans{wounds_label}:} & 8 \\
        \textbf{\trans{movement_label}:} & 4 \\
        \textbf{\trans{strength_label}:} & 2 \\
        \textbf{\trans{dexterity_label}:} & 2 \\
        \textbf{\trans{mind_label}:} & 2 \\
    \end{tabular}

    \textbf{Biss} (6)\\
    Durchschlag 1\\


    \subsection*{Zwarl}

    Eine ausgesprochen großwüchsige, fettleibige Fellbestie, die sich auf zwei Beinen fortbewegt. Sie kann bis zu 4 Schritt in die Höhe missen, ist aber auch nicht minder stark in die Breite gebaut. Ein ausgewachsener Zwarl wiegt an die 2 Fuhren.

    \begin{tabular}{@{}ll@{}}
        \textbf{\trans{wounds_label}:} & 12 \\
        \textbf{\trans{movement_label}:} & 6 \\
        \textbf{\trans{strength_label}:} & 4 \\
        \textbf{\trans{dexterity_label}:} & 2 \\
        \textbf{\trans{mind_label}:} & 4 \\
    \end{tabular}

    \textbf{Hieb} (10)\\
    Der Schlag trifft bis zu drei Gegner in einer Reichweite von 2 Metern. Er verursacht Geschockt 1.\\
    \textbf{Wälzen} (8)\\
    Wenn der Zwarl in Bedrängnis ist, wälzt er sich über seine Gegner. Er kommt dabei bis zu drei Gegner weit, wenn diese in einer Linie bis zu fünf Schritt entfern stehen. Er verursacht an jedem Gegner die Wunden dieses Wurfs.\\


    \subsection*{Gorilla}

    Ein massiger und intelligenter Primat, der für kontrollierte Stärke, soziale Dominanz und plötzliche, explosive Gewalt steht. Innerhalb seiner Gruppe meist ruhig, reagiert er bei Bedrohung mit furchteinflößender Kraft und verteidigt Revier und Gefährten mit überwältigender physischer Gewalt. Dominanzgesten – Brusttrommeln, Brüllen, Scheinangriffe – dienen oft als Warnung, doch wenn er angreift, geschieht dies schnell und brutal. Ein Gorilla jagt nicht wie ein Raubtier; er stellt sich der Gefahr und zerschmettert sie. In düsteren Auslegungen wirkt er beinahe bewusst in seinem Zorn – berechnend, beschützend und unaufhaltsam, wenn er entfesselt wird.

    \begin{tabular}{@{}ll@{}}
        \textbf{\trans{wounds_label}:} & 6 \\
        \textbf{\trans{movement_label}:} & 4 \\
        \textbf{\trans{strength_label}:} & 2 \\
        \textbf{\trans{dexterity_label}:} & 3 \\
        \textbf{\trans{mind_label}:} & 3 \\
    \end{tabular}

    \textbf{Griff} (8)\\
    Das Opfer wird festgehalten und kann sich nur mit einer Kraft Probe befreien.\\


    \subsection*{Eule}

    In den stillen Stunden bewegt sich die Eule dort, wo selbst Geräusche zu verschwinden scheinen. Ihr Flug ist nahezu lautlos, die Schwingen gleiten ohne Vorwarnung durch die Dunkelheit. Geleitet von außergewöhnlichem Gehör nimmt sie Bewegungen wahr, die anderen entgehen, und schlägt mit plötzlicher Präzision von oben zu. Sie verweilt nicht – erscheint nur kurz, um zu greifen, und verschwindet ebenso schnell wieder im Schatten. In düsteren Auslegungen wirkt sie wie eine Beobachterin der Nacht, ungesehen und doch stets präsent, ein Flüstern von Bewegung jenseits der Wahrnehmung.

    \begin{tabular}{@{}ll@{}}
        \textbf{\trans{wounds_label}:} & 2 \\
        \textbf{\trans{movement_label}:} & 8 \\
        \textbf{\trans{strength_label}:} & 1 \\
        \textbf{\trans{dexterity_label}:} & 3 \\
        \textbf{\trans{mind_label}:} & 4 \\
    \end{tabular}

    \textbf{Schnabel} (6)\\
    Durchschlag 0\\


    \subsection*{Gelber Tiger}

    Wie sein Name sagt, ist sein Fell zwischen den Streifen gelblich und heller als das der anderen Tiger. Die tropischen, von zahlreichen Flüssen durchzogenen Länder der [[ancatir|Ancatir]] und der [[ogrut|O'Grut]] sind seine Heimat.

    \begin{tabular}{@{}ll@{}}
        \textbf{\trans{wounds_label}:} & 8 \\
        \textbf{\trans{movement_label}:} & 8 \\
        \textbf{\trans{strength_label}:} & 3 \\
        \textbf{\trans{dexterity_label}:} & 5 \\
        \textbf{\trans{mind_label}:} & 2 \\
    \end{tabular}

    \textbf{Krallen} (8)\\
    Durchschlag 1, Blutend 1\\
    \textbf{Biss} (10)\\
    Durchschlag 1, Blutend 1\\


    \subsection*{Gans}

    An stillen Gewässern und auf offenen Flächen behauptet die Gans ihr Revier mit unerwarteter Entschlossenheit. Was zunächst harmlos wirkt, wird bei Annäherung schnell konfrontativ – sie zischt, breitet die Flügel aus und geht ohne Zögern nach vorn. Mit überraschender Kühnheit verteidigt sie ihr Gebiet und ihren Nachwuchs, hackt mit dem Schnabel zu und schlägt mit kräftigen Flügeln. Rückzug ist nicht ihr erster Impuls – sondern Angriff. In düsteren Auslegungen wirkt sie fast unerbittlich, eine kleine, aber wütende Wächterin, die sich unter keinen Umständen vertreiben lässt.

    \begin{tabular}{@{}ll@{}}
        \textbf{\trans{wounds_label}:} & 2 \\
        \textbf{\trans{movement_label}:} & 6 \\
        \textbf{\trans{strength_label}:} & 1 \\
        \textbf{\trans{dexterity_label}:} & 2 \\
        \textbf{\trans{mind_label}:} & 3 \\
    \end{tabular}

    \textbf{Biss} (8)\\
    Blutend 1\\


    \subsection*{Wilder Hund}

    Ein unerbittlicher Rudeljäger, der von Koordination, Ausdauer und Instinkt lebt. Wilde Hunde verlassen sich nicht auf reine Stärke – sie siegen durch Zusammenarbeit, Kommunikation und das Zermürben ihrer Beute. Über lange Strecken treiben sie ihr Ziel vor sich her, erzwingen Fehler und isolieren Schwache, bevor sie gemeinsam zuschlagen. Einzelne sind beherrschbar, doch als Rudel werden sie zu einer schnellen, disziplinierten Bedrohung. In düsteren Auslegungen wirken sie unheimlich synchron – wie ein einziger Wille, verteilt auf viele Körper.

    \begin{tabular}{@{}ll@{}}
        \textbf{\trans{wounds_label}:} & 6 \\
        \textbf{\trans{movement_label}:} & 4 \\
        \textbf{\trans{strength_label}:} & 2 \\
        \textbf{\trans{dexterity_label}:} & 2 \\
        \textbf{\trans{mind_label}:} & 2 \\
    \end{tabular}

    \textbf{Biss} (6)\\
    Blutend 1\\
    \textbf{Ruf nach dem Rudel} (6)\\
    Der wilde Hund ruft nach seinem Rudel. Nach 1W3 Kampfrunden erscheinen weitere wilde Hunde entsprechend der Erfolge dieses Wurfs.\\


    \subsection*{Grottenwurm}

    Diese bis zu drei Schritt langen Würmer leben in Gängen und Gruben vornehmlich unter der Oberfläche Tirakans. Ihre Mäuler sind mit scharfen, langen Zähnen gespickt. Zwar sind die Würmer blind, doch verfügen sie über einen hervorragenden Geruchs- und Hörsinn. [[morgalas|Morgalas]] Anführer sind dafür bekannt abgerichtete Würmer als Reittiere und im Kampf zu verwenden und diese beim Graben von Tunneln zum Einsatz zu bringen.

    \begin{tabular}{@{}ll@{}}
        \textbf{\trans{wounds_label}:} & 24 \\
        \textbf{\trans{movement_label}:} & 5 \\
        \textbf{\trans{strength_label}:} & 5 \\
        \textbf{\trans{dexterity_label}:} & 1 \\
        \textbf{\trans{mind_label}:} & 3 \\
    \end{tabular}

    \textbf{Biss} (12)\\
    Durchschlag 2, Blutend 1\\
    \textbf{Überwalzen} (14)\\
    Ist der Wurm in Bedrängnis, so wird er in einer geraden Linie auf 10 Schritt alles überwalzen, was ihm im Weg steht. Die Treffer dieses Wurfs werden auf alle überrollten Opfer aufgeteilt und verursachen Durchschlag 1.\\


    \subsection*{Riesenkrabbe}

    Über Gezeitenflächen und entlang zerklüfteter Küsten bewegt sich die Riesenkrabbe mit beunruhigender Zielstrebigkeit. Ihr schwerer Panzer hält selbst kräftige Schläge ab, während ihre gewaltigen Scheren mit zermalmender Kraft zuschnappen und Knochen wie Rüstung brechen können. Sie stürmt nicht blind voran – jede Bewegung ist bedacht, während sie sich nähert und sich zugleich mit erhobenen Scheren schützt. In ihrem Revier ist der Untergrund tückisch, und die Flucht wird durch das Gelände erschwert, das sie bestens kennt. In düsteren Auslegungen wirkt sie wie eine lebende Bastion der Küste – unnachgiebig, fremdartig und darauf ausgelegt, Unvorsichtige einem langsamen, unausweichlichen Ende entgegenzuführen.

    \begin{tabular}{@{}ll@{}}
        \textbf{\trans{wounds_label}:} & 6 \\
        \textbf{\trans{movement_label}:} & 6 \\
        \textbf{\trans{strength_label}:} & 2 \\
        \textbf{\trans{dexterity_label}:} & 2 \\
        \textbf{\trans{mind_label}:} & 1 \\
    \end{tabular}

    \textbf{Klauen} (10)\\
    \textit{Dad-a-chum? Dum-a-chum? Ded-a-chek? Did-a-chick?}\\


    \subsection*{Wolf}

    Ein räuberischer Jäger der Wildnis, der für Geduld, Instinkt und Zusammenarbeit steht. Wölfe treten selten allein auf, sondern jagen im Rudel, prüfen ihre Beute auf Schwäche und schlagen erst dann zu. Sie meiden unnötige Risiken, kreisen ihr Ziel ein, zermürben es und greifen an, wenn der Vorteil sicher ist. In düsteren Auslegungen wirken sie unnatürlich klug, unerbittlich oder beunruhigend furchtlos.

    \begin{tabular}{@{}ll@{}}
        \textbf{\trans{wounds_label}:} & 4 \\
        \textbf{\trans{movement_label}:} & 5 \\
        \textbf{\trans{strength_label}:} & 2 \\
        \textbf{\trans{dexterity_label}:} & 3 \\
        \textbf{\trans{mind_label}:} & 2 \\
    \end{tabular}

    \textbf{Biss} (7)\\
    Blutend 1\\


    \subsection*{Reitdrachen}

    Beim Reitdrachen handelt es sich um keinen echten Drachen, lediglich um eine große Flugechse mit gekörntem Kopf. Diese seltenen Tiere können, wenn sie als Jungtier gefangen werden, oder wenn man ein Gelege entdeckt und es sichern kann, zu treuen Reittieren abgerichtet werden.

    \begin{tabular}{@{}ll@{}}
        \textbf{\trans{wounds_label}:} & 8 \\
        \textbf{\trans{movement_label}:} & 8 \\
        \textbf{\trans{strength_label}:} & 4 \\
        \textbf{\trans{dexterity_label}:} & 4 \\
        \textbf{\trans{mind_label}:} & 3 \\
        \textbf{\trans{resistances_label}:} & ['Feuer'] \\
    \end{tabular}

    \textbf{Krallen} (8)\\
    Durchschlag 1, Blutend 1\\
    \textbf{Feuerstrahl} (10)\\
    Reitdrachen sind in der Lage, einen Feuerstrahl von etwa 10 Schritt Länge zu spucken. Jeder im Bereich dieses Strahls erhält die Treffer dieses Wurfs mit Durchschlag 0.\\


    \subsection*{Streifenpferd}

    Dem Pferde der nördlichen Steppen gleich, besiedeln die Streifenpferde die Steppen der [[die-stamme-der-barbaren|Barbaren]]. Die Farbe des Körpers unterteilt sich in weisse und dunkelbraune Streifen, die dem Streifenpferd sein unverwechselbares Aussehen geben.

    \begin{tabular}{@{}ll@{}}
        \textbf{\trans{wounds_label}:} & 8 \\
        \textbf{\trans{movement_label}:} & 8 \\
        \textbf{\trans{strength_label}:} & 2 \\
        \textbf{\trans{dexterity_label}:} & 3 \\
        \textbf{\trans{mind_label}:} & 2 \\
    \end{tabular}

    \textbf{Biss} (8)\\
    Durchschlag 0\\
    \textbf{Hufe} (8)\\
    Geschockt 1\\


    \subsection*{Löwe}

    Ein ausgewachsenes Löwenmännchen kann im ausgewachsenen Zustand gut 1,5 Schritt Schulterhöhe erreichen und sein Maul zieren zwei mächtige Säbelzähne.

    \begin{tabular}{@{}ll@{}}
        \textbf{\trans{wounds_label}:} & 8 \\
        \textbf{\trans{movement_label}:} & 6 \\
        \textbf{\trans{strength_label}:} & 3 \\
        \textbf{\trans{dexterity_label}:} & 3 \\
        \textbf{\trans{mind_label}:} & 2 \\
    \end{tabular}

    \textbf{Biss} (8)\\
    Durchschlag 1, Blutend 1\\
    \textbf{Krallen} (6)\\
    Durchschlag 1, Blutend 2\\


    \section{Magisch}



    \subsection*{Werwolf}

    Der Werwolf ist ein Lebewesen, halb Wolf, halb Mensch. Bei Vollmond verwandelt sich der bis dahin normale Mensch in einen Wolf, welcher jedoch bis zu zwei Schritt groß werden kann und äußerst aggressiv ist. Jeder kann ein Werwolf werden, sobald er von einem bereits infizierten Werfolf gebissen wird.

    \begin{tabular}{@{}ll@{}}
        \textbf{\trans{wounds_label}:} & 12 \\
        \textbf{\trans{movement_label}:} & 8 \\
        \textbf{\trans{strength_label}:} & 5 \\
        \textbf{\trans{dexterity_label}:} & 3 \\
        \textbf{\trans{mind_label}:} & 3 \\
    \end{tabular}

    \textbf{Biss} (8)\\
    Durchschlag 1, Blutend 1\\
    \textbf{Krallen} (8)\\
    Blutend 1\\


    \subsection*{Schillerwal}

    Der Schillerwal ist ein 50 Schritt langer massiger Wal, welcher in der Lage ist in Gefahrensituationen seine Hautfarbe der Umgebung anzupassen. Durch das Schillern wie es die Walfänger nennen kann der Wal kaum noch geortet werden. Es handelt sich hier um ein Tier, welches wohl instinktiv Magie anwenden kann.

    \begin{tabular}{@{}ll@{}}
        \textbf{\trans{wounds_label}:} & 40 \\
        \textbf{\trans{movement_label}:} & 8 \\
        \textbf{\trans{strength_label}:} & 5 \\
        \textbf{\trans{dexterity_label}:} & 3 \\
        \textbf{\trans{mind_label}:} & 4 \\
        \textbf{\trans{resistances_label}:} & ['Kälte', 'Magie', 'Feuer'] \\
    \end{tabular}

    \textbf{Rammen} (10)\\
    Der Wal vermag seine Opfer, mitunter ganze Schiffe, zu rammen. Der Wurf verursacht Durchschlag 2. Gelingt der Wurf mit mehr als 5 Erfolgen, so verschlingt der Wal sein Opfer.\\
    \textbf{Wilde Magie} (10)\\
    Der Wal wirkt einen Zauber der Elementarmagie, ein Wirken der Ginae oder des Duglaraan. Der Zauber wird intuitiv gewirkt und tritt mit einer \textbf{Stärke} gemäß den Erfolgen dieses Wurfs in Erscheinung.\\


    \subsection*{Bork}

    Diese riesenhaften Schweine sind eigentlich von gutmütiger Natur. Wird ihr Lebensraum jedoch eingeschränkt oder sie in eine Falle getrieben, so wehren sie sich. Sie sind magische Wesen, und beherrschen das magische Rammen von Gegnern.

    \begin{tabular}{@{}ll@{}}
        \textbf{\trans{wounds_label}:} & 8 \\
        \textbf{\trans{movement_label}:} & 6 \\
        \textbf{\trans{strength_label}:} & 3 \\
        \textbf{\trans{dexterity_label}:} & 3 \\
        \textbf{\trans{mind_label}:} & 3 \\
    \end{tabular}

    \textbf{Magisches Rammen} (7)\\
    Das Bork stürmt auf sein Opfer zu und verbindet Bewegung mit Angriff. Durchschlag 2\\
    \textbf{Rammen} (6)\\
    Durchschlag 1\\
    \textbf{Beissen} (5)\\
    Durchschlag 1\\


    \subsection*{Fel'war}

    Diese magischen Wesen kommen vermehrt in den Steppen der [[quitaron|Quitaron]] vor. Der Raubvogel jagt hauptsächlich kleine Nagetiere, mitunter stehen allerdings auch größere Tiere auf seinem Speiseplan.

    \begin{tabular}{@{}ll@{}}
        \textbf{\trans{wounds_label}:} & 3 \\
        \textbf{\trans{movement_label}:} & 8 \\
        \textbf{\trans{strength_label}:} & 1 \\
        \textbf{\trans{dexterity_label}:} & 4 \\
        \textbf{\trans{mind_label}:} & 2 \\
        \textbf{\trans{resistances_label}:} & ['Magie'] \\
    \end{tabular}

    \textbf{Krallen} (6)\\
    Blutend 1\\
    \textbf{Schnabel} (5)\\
    Blutend 1\\
    \textbf{Luftstoss} (5)\\
    Ein magischer Schwall Luft stößt das Opfer W6 Meter nach hinten und verursacht \textbf{Stärke} Wunden und Schock 1.\\


    \subsection*{Steinelementar}

    Das Steinelementar als Diener des [[tador-titan-des-steins|Tador]] erwächst, so es erwacht, aus einem Fels oder Stein. Steinelementare treten in der Regel als bullige humanoide Gestalten auf, welche durchaus 4 Schritt groß werden können.

    \begin{tabular}{@{}ll@{}}
        \textbf{\trans{wounds_label}:} & 30 \\
        \textbf{\trans{movement_label}:} & 4 \\
        \textbf{\trans{strength_label}:} & 6 \\
        \textbf{\trans{dexterity_label}:} & 1 \\
        \textbf{\trans{mind_label}:} & 1 \\
        \textbf{\trans{resistances_label}:} & ['Feuer', 'Gift'] \\
    \end{tabular}

    \textbf{Steinschlag} (8)\\
    Wenn das Steinelementar angreift, kann es Teile seiner eigenen Substanz auf seine Gegner schleudern. Für je drei Wunden, die es sich dabei selbst zufügt, erleiden alle Gegner in einem Kegel von 45° 1W6 Treffer. Der Kegel reicht Glaubensniveau*2 Schritte weit. Gelingt der Angriff nicht, erleidet nur das Steinelementar Schaden, die Gegner bleiben unversehrt.\\
    \textbf{Felsschlag} (6)\\
    Das Steinelementar kann einen großen Felsbrocken aus seiner eigenen Substanz werfen. Der Brocken trifft ein Ziel in bis zu 20 Schritten Entfernung und verursacht 3 Treffer pro Erfolg. Dabei fügt sich das Steinelementar selbst 2W6 Wunden zu. Misslingt der Wurf, erleidet nur das Steinelementar den Schaden und das Ziel bleibt unversehrt.\\
    \textbf{Felswuchs} (6)\\
    Das Steinelementar lässt in einem Kreis mit \textbf{Glaubensniveau} Schritt Radius spitze Felsen aus dem Boden wachsen. Dieser Kreis kann bis zu 20 Schritt vom Steinelementar entfernt sein. Beendet ein Gegner seine Kampfrunde auf den Spitzen, so erleidet er drei Treffer.

Will sich ein Gegner durch die Spitzen bewegen, muss er einen Geschickwurf machen. Gelingt ihm dieser Wurf nicht, so muss er drei Treffer hinnehmen.\\


    \subsection*{Wasserelementar}

    Das Wasserelementar ist ein elementares Wesen welches entweder unter der Gunst der Ginae oder des Duglaraan steht. Wasserelementare existieren nur auf der Elementarebene natürlich und können nur von Elementarmagier auf diese Welt gerufen werden.

    \begin{tabular}{@{}ll@{}}
        \textbf{\trans{wounds_label}:} & 14 \\
        \textbf{\trans{movement_label}:} & 8 \\
        \textbf{\trans{strength_label}:} & 2 \\
        \textbf{\trans{dexterity_label}:} & 5 \\
        \textbf{\trans{mind_label}:} & 2 \\
        \textbf{\trans{resistances_label}:} & ['Wasser', 'Magie', 'Feuer', 'Physikalischer Schaden'] \\
    \end{tabular}

    \textbf{Spiegelbild} (8)\\
    Ist der Wurf auf Spiegelbild gelungen  so erzeugt das Wasserelementar einen genauen Klon von sich, welcher alle seine Werte erhält. Die Fähigkeit kostet 1W6 Arkana, welche VOR dem Klonen abgezogen werden.\\
    \textbf{Schlag} (8)\\
    Durchschlag 1\\
    \textbf{Wuchs} (6)\\
    Das Wasserelementar wächst und verdoppelt sowohl seine maximalen, als auch seine aktuellen Wunden. Sobald die maximalen Wunden 100 erreichen wächst es nicht weiter.\\
    \textbf{Greifen} (8)\\
    Greift das Wasserelementar ein Opfer, so zieht es dieses in seinen Körper. Das Opfer bekommt in jeder seiner Kampfrunden zwei Wunden Erstickungsschaden, kann sich jedoch mit einer Kraft oder Akrobatkprobe nach dem Schaden befreien.\\


    \subsection*{Kompostfee}

    Die Kompostfee ist ein sonderbares, magisches Wesen. Sie lebte einst im Kompost der Hexe Mare, im Mittelalter auf der Erde, unweit der Stadt Aquisgrani. Mare und sie verbinden ein besonderes Band, welches auch auf magische Weise geprägt ist.

    \begin{tabular}{@{}ll@{}}
        \textbf{\trans{wounds_label}:} & 2 \\
        \textbf{\trans{movement_label}:} & 8 \\
        \textbf{\trans{strength_label}:} & 2 \\
        \textbf{\trans{dexterity_label}:} & 5 \\
        \textbf{\trans{mind_label}:} & 3 \\
        \textbf{\trans{resistances_label}:} & ['Magie'] \\
    \end{tabular}

    \textbf{Energiekugel} (2)\\
    Durchschlag 2, Geschockt 1\\


    \subsection*{Sethlarn}

    Die Sethlarn sind eine magische Art von drachenartigen Wesen. Ihre Gestalt ist die von [[die-drachen|Drachen]], jedoch ist ihre ledrige Haut grob und schwarz. Ihre Augen schimmern in einem fahlen Gelb. Sie erreichen eine Größe von drei Schritt, und eine Flügelspannweite von fünf Schritt.

    \begin{tabular}{@{}ll@{}}
        \textbf{\trans{wounds_label}:} & 18 \\
        \textbf{\trans{movement_label}:} & 10 \\
        \textbf{\trans{strength_label}:} & 3 \\
        \textbf{\trans{dexterity_label}:} & 4 \\
        \textbf{\trans{mind_label}:} & 3 \\
        \textbf{\trans{resistances_label}:} & ['Magie', 'Feuer', 'Schock'] \\
    \end{tabular}

    \textbf{Krallen} (10)\\
    Durchschlag 1, Bleeding 1\\
    \textbf{Dunkle Magie} (8)\\
    Die Sethlarn beherrschen die Zauber der Schwarzen und der Echsenmagie, die sie intuitiv einsetzen. Gelingt dieser Wurf, so tritt der Effekt eines Zaubers mit \textbf{Stärke} entsprechend den Erfolgen in Kraft.\\


    \subsection*{Imps}

    Die Imps sind kleine 30cm hohe humanoide Wesen von dürrer knorriger Gestalt und kleinen roten funkelnden Augen. Ihre runzelige Haut ist von kleinen Knötchen bedeckt und von grauer bis brauner Farbe. Auch wenn es so scheint, sind nicht mit den Bolden verwandt.

    \begin{tabular}{@{}ll@{}}
        \textbf{\trans{wounds_label}:} & 4 \\
        \textbf{\trans{movement_label}:} & 3 \\
        \textbf{\trans{strength_label}:} & 2 \\
        \textbf{\trans{dexterity_label}:} & 3 \\
        \textbf{\trans{mind_label}:} & 4 \\
    \end{tabular}

    \textbf{Krallen} (5)\\
    Blutend 1\\
    \textbf{Elementarmagie} (10)\\
    Der Imp spricht einen Zauber der Elementarmagie, der seinem Element zugetan ist. Die Erfolge dieses Wurfs ergeben die \textbf{Stärke} des Zaubers.\\


    \subsection*{Ganark}

    Ganark ist ein Wesen, welches in verzauberten Wäldern anzutreffen ist. Vom Habitus her ähnelt er einem Kobold, es bereitet ihm Unterhaltung, Fremde zu verwirren.

    \begin{tabular}{@{}ll@{}}
        \textbf{\trans{wounds_label}:} & 4 \\
        \textbf{\trans{movement_label}:} & 6 \\
        \textbf{\trans{strength_label}:} & 2 \\
        \textbf{\trans{dexterity_label}:} & 3 \\
        \textbf{\trans{mind_label}:} & 5 \\
        \textbf{\trans{resistances_label}:} & ['Magie', 'Aktionen stehlen'] \\
    \end{tabular}

    \textbf{Verwirren} (8)\\
    Das Opfer ist für 1W3 Kampfrunden verwirrt. Der Mindestwurf aller Würfe ist um die Anzahl der Erfolge angehoben, die Ganark bei diesem Verwirren Wurf erzielt.\\
    \textbf{Kratzen} (4)\\
    Ganark kratzt das Opfer (Durchschlag 1)\\
    \textbf{Verschwimmen} (6)\\
    Ganark bricht das Licht um sich und lässt sein Form verschwimmen. Für \textit{Zauberstärke} Runden sind Angriffe gegen ihn erschwert (Mindestwurf +1).\\


    \subsection*{Bolde}

    Nur selten sieht man Bolde in den Wäldern und Städten Tirakans. Begabte sprechen ihnen einen Ursprung in den Feenwelten zu, während das normale Volk Bolde als Unglücksbringer und Scharlatane verächtet. Ob jemals geklärt wird, wo der Ursprung der Bolde liegt, mag bezweifelt werden. Tatsächlich gibt es in der Feenwelt viele Bolde.

    \begin{tabular}{@{}ll@{}}
        \textbf{\trans{wounds_label}:} & 4 \\
        \textbf{\trans{movement_label}:} & 3 \\
        \textbf{\trans{strength_label}:} & 1 \\
        \textbf{\trans{dexterity_label}:} & 3 \\
        \textbf{\trans{mind_label}:} & 5 \\
        \textbf{\trans{resistances_label}:} & ['Magie'] \\
    \end{tabular}

    \textbf{Wilde Magie} (8)\\
    Der Bold wirkt einen Zauber der Elementarmagie, der Zauberei oder des Schamanismus. Die Erfolge dieses Wurfs stellen die \textbf{Stärke} des Zaubers dar.\\


    \subsection*{Nymphen}

    Nymphen sind Diener der [[die-titanen|Titanin]] [[ginae-titanin-des-wassers|Ginae]], eineinhalb Schritt hohe Frauengestalten, die Merkmale von Wasserbewohnern an ihrem Leibe tragen. Ihre liebste Beschäftigung ist das Spiel mit der Harfe oder der Seemuschelflöte und das Verabreichen von feuchten Küssen.

    \begin{tabular}{@{}ll@{}}
        \textbf{\trans{wounds_label}:} & 6 \\
        \textbf{\trans{movement_label}:} & 4 \\
        \textbf{\trans{strength_label}:} & 2 \\
        \textbf{\trans{dexterity_label}:} & 2 \\
        \textbf{\trans{mind_label}:} & 3 \\
        \textbf{\trans{resistances_label}:} & ['Magie'] \\
    \end{tabular}

    \textbf{Ginaes Wirken} (6)\\
    Nymphen können die Titanin Ginae anrufen und ihr Wirken erbitten.\\
    \textbf{Wassermagie} (10)\\
    Nymphen können Zauber der Elementarmagie und Zauberei sprechen, die in Verbindung mit dem Wasser stehen. Die Erfolge dieses Wurfs geben die \textbf{Stärke} des Zaubers an.\\


    \subsection*{Schleimhüpfer}

    Die Gestalt eines Schleimhüpfers ist am ehesten mit der eines kleinen Boldes, welcher sich gedrungen und dem Boden nahe fortbewegt. Der Name rührt von der Art der Fortbewegung her, der Schleimhüpfer bewegt sich mit einem hüpfenden, fast taumelnden Gang fort. Sein dunkler Körper ist bedeckt mit einem Schleim, den der Schleimhüpfer selbst erzeugt. Der Schleim brennt auf der Haut, und verursacht eine deutliche Rötung.

    \begin{tabular}{@{}ll@{}}
        \textbf{\trans{wounds_label}:} & 4 \\
        \textbf{\trans{movement_label}:} & 3 \\
        \textbf{\trans{strength_label}:} & 2 \\
        \textbf{\trans{dexterity_label}:} & 2 \\
        \textbf{\trans{mind_label}:} & 3 \\
    \end{tabular}

    \textbf{Brennender Schleim} (6)\\
    Wird man von dem Schleim des Schleimhüpfers benetzt, so verursacht der Schleim eine Wunde pro Kampfrunde, bis er abgewischt wird. Dabei brennt die benetzte Stelle fürchterlich.\\
    \textbf{Ersticken} (4)\\
    Kommt der Schleimhüpfer in Bedrängnis, so versucht er, dem Angreifer ins Gesicht zu springen. Gelingt ihm dies (ein Erfolg auf dieser Aktion), so nimmt das Opfer zwei Wunden pro Kampfrunde hin, bis es sich mit einer Kraftprobe von dem Schleimhüpfer befreit hat.\\


    \subsection*{Riesen}

    Sowohl in den nördlichen Wäldern nahe des [[tal-des-vergessens|Tal des Vergessens]], als auch im Süden unweit der Gebiete der [[xordai|Xordai]] gibt es Riesen. Diese gigantischen, menschenartigen Wesen leben ihre eigene Kultur, haben ihre eigene Sprache, und halten sich von den [[menschen|Menschen]] fern.

    \begin{tabular}{@{}ll@{}}
        \textbf{\trans{wounds_label}:} & 25 \\
        \textbf{\trans{movement_label}:} & 8 \\
        \textbf{\trans{strength_label}:} & 6 \\
        \textbf{\trans{dexterity_label}:} & 2 \\
        \textbf{\trans{mind_label}:} & 2 \\
    \end{tabular}

    \textbf{Hieb} (12)\\
    Schock 1\\


    \subsection*{Ventrikulum}

    Eines der wohl seltsamsten Wesen Tirakans ist das Ventrikulum. Wenngleich Beobachter der vergangenen Jahrhunderte es als untotes Wesen klassifizierten muss es doch eher als rein magische Entität gelten. Trotzdem besteht das Ventrikulum aus den toten Überresten eines [[menschen|Menschen]], welche magisch zum Leben erweckt wurden.

    \begin{tabular}{@{}ll@{}}
        \textbf{\trans{wounds_label}:} & 15 \\
        \textbf{\trans{movement_label}:} & 6 \\
        \textbf{\trans{strength_label}:} & 1 \\
        \textbf{\trans{dexterity_label}:} & 3 \\
        \textbf{\trans{mind_label}:} & 5 \\
        \textbf{\trans{resistances_label}:} & ['Kälte', 'Magie', 'Gift', 'Schock'] \\
    \end{tabular}

    \textbf{Telekinese} (8)\\
    Mittels Konzentration vermag das Ventrikulum Objekte im Umkreis von 10 Schritt bis zu einem Gewicht von 50 Stein zu bewegen. Für je ein Arkana pro 3 Stein kann das Gewicht diese 50 Stein auch überschreiten.

Der Gegenstand kann mit beliebiger Geschwindigkeit bewegt werden, so dass auch Gegenstände mit großer Geschwindigkeit gegen ein Opfer geschleudert werden können. Der Schaden hierbei wird vom Spielleiter festgelegt und richtet sich nach Art und Geschwindigkeit des Gegenstands. Es können auch Lebende Wesen bewegt werden.\\
    \textbf{Geistkontrolle} (6)\\
    Das Ventrikulum kontrolliert den Geist eines intelligenten Wesens. Hierbei kann das Ventrikulum die Handlungen des Opfers komplett kontrollieren, muss jedoch jede Kampfrunde eine seiner Aktionen dafür aufwenden die Kontrolle aufrecht zu erhalten.\\
    \textbf{Dyspnoe} (6)\\
    Das Ventrikulum vermag es ein Opfer magisch zu würgen. Um das Würgen aufrecht zu erhalten muss jede Kampfrunde eine Aktion aufgewendet werden. Das gewürgte Opfer nimmt in der ersten Kampfrunde eine Wunde. Der Schaden verdoppelt sich jede Kampfrunde.\\
    \textbf{Sprung} (8)\\
    Innerhalb einer Aktion vermag das Ventrikulum zu einer beliebigen sichtbaren Position im Umkreis von 10 Schritt zu springen.\\
    \textbf{Blitzgewitter} (12)\\
    Aus den Nervenenden des Ventrikulums strömen Blitze, welche in einem Kegel von 90 Grad bei einer Reichweite von 10 Schritt an allen darin befindlichen Lebewesen Treffer entsprechend der Erfolge dieses Wurfs pro Kampfrunde macht. Um das Blitzgewitter aufrecht zu erhalten muss das Ventrikulum eine Aktion pro Runde aufwenden.\\


    \subsection*{Chronars Segen}

    Dieser als [[chronar-gott-der-zeit|Chronars]] Segen bekannter Falter kommt nur sehr selten auf Tirakan vor. Man sagt ihm nach er lande nachts auf dem Haupt von schwer Kranken oder Verwundeten, heilt diese oder begleitet sie schmerzlindernd in das Reich der Toten.

    \begin{tabular}{@{}ll@{}}
        \textbf{\trans{wounds_label}:} & 1 \\
        \textbf{\trans{movement_label}:} & 2 \\
        \textbf{\trans{strength_label}:} & 1 \\
        \textbf{\trans{dexterity_label}:} & 2 \\
        \textbf{\trans{mind_label}:} & 4 \\
    \end{tabular}



    \subsection*{Babas Haus}

    Baba baute eine Hütte auf dem verrottenden Baumstumpf eines riesigen Baumes, der vor langer Zeit gefällt wurde. Erst nachdem sie einen magischen Edelstein in die Hütte eingebettet hatte, wurde das Ganze mit einem Anschein von Leben erfüllt. Wenn sie es will, reißt die Hütte ihre riesigen Wurzeln aus der Erde und schlurft wie ein spinnenartiges Ungetüm umher, wobei sie bei jedem Schritt den Boden erschüttert. Die Hütte greift mit ihren um sich schlagenden und stampfenden Wurzeln an. Mit seinen Wurzeln kann er auch große Steine ​​schleudern.

    \begin{tabular}{@{}ll@{}}
        \textbf{\trans{wounds_label}:} & 20 \\
        \textbf{\trans{movement_label}:} & 4 \\
        \textbf{\trans{strength_label}:} & 5 \\
        \textbf{\trans{dexterity_label}:} & 1 \\
        \textbf{\trans{mind_label}:} & 4 \\
        \textbf{\trans{resistances_label}:} & ['Feuer', 'Aktionen stehlen'] \\
    \end{tabular}

    \textbf{Wurzelangriff} (10)\\
    Angriff der Wurzel\\
    \textbf{Steinwurf} (8)\\
    Steinwurf\\


    \subsection*{Basilisk}

    Von allen uns bekannten Wesen und Geschöpfen ist keines gefährlicher und todbringender als der Basilisk - anders auch genannt der König der Schlangen. Jenes Wesen wird aus einem Hühnerei geboren, dass von einer Krähe ausgebrütet wird. Das Gift des Basiliken ist tödlich.

    \begin{tabular}{@{}ll@{}}
        \textbf{\trans{wounds_label}:} & 12 \\
        \textbf{\trans{movement_label}:} & 4 \\
        \textbf{\trans{strength_label}:} & 3 \\
        \textbf{\trans{dexterity_label}:} & 4 \\
        \textbf{\trans{mind_label}:} & 5 \\
        \textbf{\trans{resistances_label}:} & ['Magie', 'Feuer', 'Schock', 'Aktionen stehlen'] \\
    \end{tabular}

    \textbf{Biss} (10)\\
    Durchschlag 1, Vergiften 6\\
    \textbf{Tödlicher Blick} (10)\\
    Der Blick des Basilisken ist gleichermaßen paralysierend als auch tödlich. Trifft der Blick jemanden direkt verursacht er 20 Treffer. Trifft er jemanden über eine Reflektion so verursacht er 2W6 Treffer und paralysiert das Opfer auf unbestimmte Zeit.\\


    \subsection*{Einhorn}

    Das Einhorn (von [[silkanda|silk.]] licorne - das Reine) ist ein Tier, welches dem Pferd in Gestalt und Größe ähnelt, jedoch schlanker gebaut ist und auf seiner Stirn ein Horn trägt.

    \begin{tabular}{@{}ll@{}}
        \textbf{\trans{wounds_label}:} & 14 \\
        \textbf{\trans{movement_label}:} & 8 \\
        \textbf{\trans{strength_label}:} & 3 \\
        \textbf{\trans{dexterity_label}:} & 3 \\
        \textbf{\trans{mind_label}:} & 8 \\
        \textbf{\trans{resistances_label}:} & ['Magie'] \\
    \end{tabular}

    \textbf{Wilde Magie} (12)\\
    Das Einhorn kann Zauber der Elementarmagie, der Weissen Magie und des Schamanismus wirken. Die Erfolge dieses Wurfs werden die \textbf{Stäke} des Zaubers.\\


    \subsection*{Gemeine Flugechse}

    Die gemeine Flugechse hat zwei Beine und zwei ledrige Flügel. Auf dem Boden müßen sie sich mit ihren Flügeln abstützen, in der Luft aber sind diese schlanken Echsen flinkge und schnelle Flieger. Ihre mit kleinen grünen Schuppen bedeckte Haut, bietet ihr einen recht guten Schutz, und besonders kräftigen Flugechsen kann man auch noch eine leichte metallverstärkte Lederrüstung anlegen.

    \begin{tabular}{@{}ll@{}}
        \textbf{\trans{wounds_label}:} & 10 \\
        \textbf{\trans{movement_label}:} & 10 \\
        \textbf{\trans{strength_label}:} & 3 \\
        \textbf{\trans{dexterity_label}:} & 2 \\
        \textbf{\trans{mind_label}:} & 4 \\
    \end{tabular}

    \textbf{Biss} (8)\\
    Durchschlag 1, Blutend 1\\
    \textbf{Krallen} (8)\\
    Durchschlag 2\\


    \subsection*{Schattenkriecher}

    Erschaffen vom Schwarzmagier Ortan sind die Schattenkriecher humanoide Wesen. Sie leben in den nördlichen Wäldern Asgorans, wo sie künstlich durch Magie erschaffen wurden.

    \begin{tabular}{@{}ll@{}}
        \textbf{\trans{wounds_label}:} & 8 \\
        \textbf{\trans{movement_label}:} & 5 \\
        \textbf{\trans{strength_label}:} & 2 \\
        \textbf{\trans{dexterity_label}:} & 4 \\
        \textbf{\trans{mind_label}:} & 2 \\
        \textbf{\trans{resistances_label}:} & ['Magie', 'Licht'] \\
    \end{tabular}

    \textbf{Krallen} (6)\\
    Durchschlag 1\\
    \textbf{Schattenwolke} (8)\\
    Der Schattenkriecher spuckt eine Wolke aus reinem Schatten, die alles Licht in einer Sphäre mit einem Radius entsprechend der Erfolge des Wurfs verschlingt. Diese vollkommene Finsternis erhöht den Mindestwurf aller Charaktere außer dem Schattenkriecher um 4.

Die Sphäre bleibt für 2 Kampfrunden bestehen und richtet keinen Schaden an. Nachtsicht verhindert den Malus auf den Mindestwurf nicht.\\


    \subsection*{Baumplage}

    Erwachte Pflanzen, ausgestattet mit den Kräften der Intelligenz und der Beweglichkeit, überfallen die von der Finsternis verdorbenen Länder. Indem sie die Finsternis aus dem Boden trinken, erfüllen sie den Willen eines uralten Bösen und versuchen, dieses Böse zu verbreiten, wo immer sie können.

    \begin{tabular}{@{}ll@{}}
        \textbf{\trans{wounds_label}:} & 6 \\
        \textbf{\trans{movement_label}:} & 2 \\
        \textbf{\trans{strength_label}:} & 4 \\
        \textbf{\trans{dexterity_label}:} & 1 \\
        \textbf{\trans{mind_label}:} & 2 \\
        \textbf{\trans{resistances_label}:} & ['Wasser', 'Feuer'] \\
    \end{tabular}

    \textbf{Wurzelpeitsche} (7)\\
    Vergiftet 1\\


    \subsection*{Beseelte Statue}

    Manchmal wird ein Geist oder eine magische Präsenz in eine Statue gebannt. Diese Statuen sind aus Stein, von ihrem Schöpfer gemeißelt. Aber der Geist der Vergangenheit lebt in ihnen und erweckt sie zu geisterhaftem Leben. Die Seele ihres Vorbildes verbindet sich mit dem steinernen Körper und wird lebendig.

    \begin{tabular}{@{}ll@{}}
        \textbf{\trans{wounds_label}:} & 12 \\
        \textbf{\trans{movement_label}:} & 2 \\
        \textbf{\trans{strength_label}:} & 5 \\
        \textbf{\trans{dexterity_label}:} & 1 \\
        \textbf{\trans{mind_label}:} & 2 \\
        \textbf{\trans{resistances_label}:} & ['Magie', 'Feuer', 'Gift', 'Schock'] \\
    \end{tabular}

    \textbf{Temporaler Griff} (5)\\
    Die Statue greift nach dem Opfer. Ist der Wurf gelungen, legt sie ihre kalten Hände auf den Körper des Opfers und blickt ihm in die Augen.

In diesem Moment scheint das Opfer für den Betrachter zu verschwinden. Es ist jedoch nicht tot, sondern in eine Zeit vor 200 Jahren zurückversetzt. 

Einmal pro Stunde darf das Opfer auf seinen Widerstand werfen. Gelingt der Wurf, so wird das Opfer in die heutige Zeit zurückversetzt. Misslingt der Wurf, bleibt das Opfer in der Zeit.\\
    \textbf{Blinzeln} (8)\\
    Die beseelte Statue teleportiert sich an den gewünschten Zielort. Dieser kann bis zu \textbf{Stärke} Meter entfernt sein. Die Aktion ist nur möglich, wenn niemand auf die Statue schaut.\\


    \subsection*{Schillerwal (Leviatan)}

    Der Schillerwal ähnelt in seiner Form einem irdischen Blauwal, ist jedoch schlanker und besitzt längere, fast flügelartige Seitenflossen. Das Besondere ist seine Haut: Sie ist nicht grau, sondern besitzt eine perlmuttartige, semitransparente Oberfläche.

Je nach Lichteinfall und magischer Sättigung der Umgebung bricht seine Haut das Licht in alle Farben des Spektrums. Daher der Name. Wenn ein Schillerwal an der Oberfläche bricht, sieht es aus, als würde ein lebender Regenbogen aus dem Wasser steigen.
Gelehrte glauben, dass diese Wale sich nicht nur von Krill ernähren, sondern auch das Licht des Mondes und der Sterne absorbieren, wenn sie nachts an die Oberfläche kommen.

Schillerwale meiden die flachen Küstengewässer. Sie ziehen durch die tiefen Ozeane, weit abseits der Routen der Handelsschiffe.

Man sichtet sie am häufigsten im Südmeer, in den kalten Strömungen fernab der Hitze des Dschungels, oder in den mystischen Gewässern rund um verlassene Inselarchive.

Sie reisen in kleinen Familienverbänden (Pods). Es heißt, ihr Gesang kann Stürme beruhigen oder Wahnsinn bei jenen auslösen, die ihn zu lange hören.

Da die Jagd extrem gefährlich ist (Schillerwale wehren sich selten, werden aber oft von See-Elementaren oder Meerjungfrauen beschützt) und die Tiere selten sind, ist der Preis enorm.

Die Ancatir betrachten die Jagd auf Schillerwale als Frevel. Sie nutzen nur Tran, der von natürlich gestrandeten Walen stammt (Geschenk der Gezeiten). Alchemisten, die blutigen Tran verwenden, werden in elfischen Enklaven oft der Stadt verwiesen.

Der Schillerwal gehört zu den faszinierendsten und friedliebendsten Giganten der Meere um Tirakan. Er ist nicht nur ein Tier, sondern eine lebende Anomalie, die eng mit den magischen Strömungen der Ozeane verbunden ist.

Schillerwale greifen niemals zuerst an. Wird er getötet, erleiden alle Beteiligten einen Fluch des Meeres (+2 Erschöpfung pro Tag auf See).

Wir sahen ihn bei Neumond. Er leuchtete unter dem Kiel wie eine versunkene Stadt. Als wir die Harpunen warfen, schrie das Biest nicht. Es fing an zu singen. Ein Ton, so tief, dass das Holz meines Schiffs splitterte und zwei meiner Männer einfach ins Wasser sprangen, mit einem Lächeln im Gesicht. Wir haben ihn erlegt, ja. Aber das Öl... es brennt in den Lampen meiner Kajüte, und ich schwöre, in den Schatten sehe ich die Gesichter derer, die gesprungen sind.
Aus dem Logbuch des Walfängers 'Haken-Ulf'

    \begin{tabular}{@{}ll@{}}
        \textbf{\trans{wounds_label}:} & 80 \\
        \textbf{\trans{movement_label}:} & 20 \\
        \textbf{\trans{strength_label}:} & 16 \\
        \textbf{\trans{dexterity_label}:} & 12 \\
        \textbf{\trans{mind_label}:} & 11 \\
    \end{tabular}

    \textbf{Sonarstoß} (16)\\
    Der Wal stößt einen extrem lauten, gerichteten Schallimpuls aus.

Die Ziele müssen eine Resistenz-Probe (WI) bestehen. Bei Misserfolg sind sie für 1W6 Runden betäubt und erleiden einen Malus von +4 auf alle Proben.\\
    \textbf{Gewaltiger Flossenschlag} (14)\\
    Ein Schlag mit der massiven Schwanzflosse, der das Wasser peitscht.

3W6+16 Schaden. Trifft alle Ziele in einem 5-Meter-Radius hinter dem Wal.\\
    \textbf{Irisierendes Leuchten} (14)\\
    Der Wal lässt seine Haut in einem blendenden Farbspektrum aufblitzen.

Alle Angreifer im Nahbereich müssen eine Wahrnehmungs-Probe bestehen. Bei Misserfolg sind sie für die nächste Runde geblendet und können keine Angriffe ausführen.\\
    \textbf{Abtauchen} (12)\\
    Der Wal zieht sich in die Tiefe zurück.

Erhöht den RS für eine Runde um weitere +4 Schilde, während er sich aus dem Kampfbereich entfernt.\\


    \section{Paranormal}



    \subsection*{Phantomkrieger}

    Ein Phantomkrieger ist der untote Geist eines Soldaten oder eines anderen fähigen Waffenträgers (d. h. eines Berufskillers), der in einem gewaltsamen Konflikt oder vor Erledigung einer zugewiesenen Aufgabe gestorben ist. Obwohl er körperlos ist, trägt ein Phantomkrieger eine geisterhafte Rüstung, einen Schild und eine Waffe. Die Erinnerung eines Phantomkriegers an die Tage vor seinem Tod ist bestenfalls verschwommen.

    \begin{tabular}{@{}ll@{}}
        \textbf{\trans{wounds_label}:} & 10 \\
        \textbf{\trans{movement_label}:} & 4 \\
        \textbf{\trans{strength_label}:} & 3 \\
        \textbf{\trans{dexterity_label}:} & 3 \\
        \textbf{\trans{mind_label}:} & 2 \\
        \textbf{\trans{resistances_label}:} & ['Magie', 'Physikalischer Schaden'] \\
    \end{tabular}

    \textbf{Schwerthieb} (10)\\
    Durchschlag 1\\

\end{multicols}